To make the model concrete, the paper should analyze both digital-heavy and embodied or hybrid scenarios. A useful digital-heavy baseline already exists in Boreal's deterministic fixture set: a support-triage automation request in which one buyer opens a request, one supply responds, one commitment is accepted, fulfillment proceeds, artifacts are delivered, and settlement closes on the same durable thread. This scenario is valuable because it shows that the model handles an AI-heavy or operations-heavy request without needing to invent secondary root objects.

However, the embodied gap appears more clearly in hybrid scenarios. Consider an onsite property inspection that also requires a structured written report. A tool-first system can readily generate the report template and checklist, but it may omit the actual visit, the access coordination, the time window, and the evidence obligations. In a request-rooted model, those conditions remain attached to the request and shape both supply selection and completion validation. The digital report is a supporting artifact, not a substitute for the inspection itself.

Another useful scenario is a local inventory audit with photo evidence and reconciliation. Here the fulfillment path mixes physical presence with digital follow-up. A field-capable human may count or inspect items onsite, while an agent or tool normalizes records, produces summaries, or highlights discrepancies. The key point is that the digital support lane should remain attached to, and subordinate to, the embodied execution lane rather than silently replacing it.

Case studies of this kind help distinguish between two superficially similar systems: one that generates convincing documents about work and one that keeps the work thread honest enough to route, verify, and settle the actual underlying task. That distinction is central to the paper's claim.
